\chapter{Saját képadatbázis építése}

A mesterséges intelligencia alapú képosztályozás hatékony működésének alapfeltétele a megfelelő minőségű és kellően reprezentatív tanítóadatbázis. Jelen dolgozat egyik legfontosabb és legidőigényesebb eleme volt egy ilyen képadatbázis összeállítása, amely nemcsak mennyiségében, hanem tartalmában és szerkezetében is jól illeszkedik a konvolúciós neurális hálózatok igényeihez. Mivel rendelkezésemre állt egy saját fotóarchívum, amely két évtizednyi, gondosan rendszerezett és archivált felvételt tartalmaz, adódott a lehetőség, hogy egyéni képtáramból építsek fel egy teljes értékű osztályozási adatbázist.

A kiindulási anyag döntő többsége NEF (Nikon Electronic Format) típusú RAW formátumban volt elérhető, amely minden színcsatorna és világossági komponens részletes információit tartalmazza. Ez a formátum ideális alapot nyújt a professzionális feldolgozáshoz, ugyanakkor szükségessé tette egy egyedi konverziós pipeline megtervezését és megvalósítását. A cél az volt, hogy a képek átalakítása során a lehető legkevesebb információ vesszen el, ugyanakkor teljesüljön minden olyan követelmény, amely a neurális hálózatok betanításához szükséges: fix képméret, egységes képarány, szabványos formátum, egyértelmű osztályozás, ismétlődésmentes elrendezés és metaadatok megőrzése.

A teljes fotóarchívum mérete meghaladta a 2 terabájtot, így a képek feldolgozása szakaszosan, 10 GB-os csomagokban történt, a saját gépem SSD tárhelyének igénybevételével. Ez nemcsak a feldolgozási sebességet növelte, hanem csökkentette az archiváló háttértárolók elhasználódásának kockázatát is. A képfeldolgozó pipeline automatikusan hajtotta végre a következő lépéseket: a NEF fájlok EXIF metaadatait megtartó konverzióját JPEG formátumba (80-as tömörítéssel), méretezést 224 képpont szélességre, középre vágást 1:1 képarány eléréséhez, MD5-alapú fájlnév-képzést az ismétlődések kiszűrésére, valamint osztályozási és szűrési lépéseket. A kezdeti automatikus osztályozást többszörös kézi ellenőrzés és újraosztályozás követte, így fokozatosan alakult ki az adatbázis végleges struktúrája.

A kategóriák számát és tartalmát az osztályok kiegyensúlyozottságának elve, valamint a képanyag tartalmi homogenitása alapján határoztam meg. Az alacsony elemszámú osztályokat átcsoportosítottam, a túl nagy kategóriákat finomabb alkategóriákra bontottam, az irreleváns vagy túl heterogén osztályokat pedig elhagytam. A végső adatbázis több mint 180 000 képet tartalmaz, gondosan strukturált mappaszerkezetben, amely a tanító, validációs és tesztadatok 80-10-10 százalékos megosztását is tartalmazza.

A képadatbázis építése közben számos saját fejlesztésű segédprogram készült, amelyek automatizálták és nyomon követhetővé tették az egyes lépéseket. Ezek a programok nemcsak a képfeldolgozásért, hanem a statisztikai elemzésekért, vizualizációkért és a kategóriaeloszlás értékeléséért is felelősek. A fejezet hátralévő részében részletesen bemutatom az egyes lépéseket, a megvalósított pipeline-t, a feldolgozási logikát, valamint az e szakaszban keletkezett eredményeket és tapasztalatokat.

%ITT TARTOK----------------------

\section{A forrásképek}
A képadatbázis építése közben számos saját fejlesztésű segédprogram készült, amelyek automatizálták és nyomon követhetővé tették az egyes lépéseket. Ezek a programok nemcsak a képfeldolgozásért, hanem a statisztikai elemzésekért, vizualizációkért és a kategóriaeloszlás értékeléséért is felelősek. A fejezet hátralévő részében részletesen bemutatom az egyes lépéseket, a megvalósított pipeline-t, a feldolgozási logikát, valamint az e szakaszban keletkezett eredményeket és tapasztalatokat.

        \subsection{Motiváció}
        A képadatbázis építése közben számos saját fejlesztésű segédprogram készült, amelyek automatizálták és nyomon követhetővé tették az egyes lépéseket. Ezek a programok nemcsak a képfeldolgozásért, hanem a statisztikai elemzésekért, vizualizációkért és a kategóriaeloszlás értékeléséért is felelősek. A fejezet hátralévő részében részletesen bemutatom az egyes lépéseket, a megvalósított pipeline-t, a feldolgozási logikát, valamint az e szakaszban keletkezett eredményeket és tapasztalatokat.

        \subsection{Források}
        A képadatbázis építése közben számos saját fejlesztésű segédprogram készült, amelyek automatizálták és nyomon követhetővé tették az egyes lépéseket. Ezek a programok nemcsak a képfeldolgozásért, hanem a statisztikai elemzésekért, vizualizációkért és a kategóriaeloszlás értékeléséért is felelősek. A fejezet hátralévő részében részletesen bemutatom az egyes lépéseket, a megvalósított pipeline-t, a feldolgozási logikát, valamint az e szakaszban keletkezett eredményeket és tapasztalatokat.

    \section{Forrás adathalmaz elemzése}

        \subsection{Darabszámok}

        \subsection{Méretek}

        \subsection{Kategóriák}

    \section{Adatbázis szerkezete}

        \subsection{Képméret, formátum}

        \subsection{Osztályozás}

        \subsection{Kategóriák}

        \subsection{Mappaszerkezet}

    \section{Előfeldolgozás}

        \subsection{Kötegelt adatgyűjtés}

        \subsection{Méret és teljesítmény koráltok}

    \section{Saját feldolgozó futószalag}

        \subsection{Tulajdonságok}

        \subsection{Szkript futószalag}

            \subsubsection{Konvertálás}

            \subsubsection{Méretezés}

            \subsubsection{Vágás}

            \subsubsection{Átnevezés}

            \subsubsection{Ismétlődés szűrése}

            \subsubsection{Kategorizálás}

            \subsubsection{Csomagolás, véglegesítés}

        \subsection{Performancia}

        \subsection{Automata osztályozás}

        \subsection{Places365}

        \subsection{ImageNet1K}

        \subsection{Kézi osztályozás}

        \subsection{Véglegesítés}

    \section{Utófeldolgozás}

        \subsection{Kézi Adattisztítás, szűrés, törlés}

        \subsection{Alkategóriák elminiálása}

        \subsection{300 kép alatti kategóriák megszüntetése}

        \subsection{A rendezés finomítása, kézi és automata}

        \subsection{Arcfelismerés}
        
        \subsection{Nagy kategóriák szétosztása}

        \subsection{Üres kategóriák megszüntetése}

    \section{Kész adatbázis elemzése}

    \section{Statisztikák}

        \subsection{Darabszámok, méretek}

        \subsection{Mappavizsgáló szkript}

        \subsection{Kategóriák, statisztikák, méretek}

        \subsection{EXIF Statisztika szkript}

            \subsubsection{Tulajdonságok}

        \subsection{EXIF statisztika}

    \section{Az elkészült adatbázis értékelése}
        